\chapter{泛函分析与算子代数}

\label{chap:03}

如果说第 1 章解决的是“在什么拓扑环境里谈极限”，第 2 章解决的是“在什么可测环境里谈积分与随机性”，那么本章要回答的就是：一旦真正进入无穷维优化，计算和证明究竟在哪些线性空间里发生。有限维凸优化把大多数结构都藏在矩阵与欧氏内积后面；到了无穷维，范数、对偶、弱拓扑、伴随和谱必须被分别说清，否则存在性、对偶性与算法收敛就会停留在口号层面。

本章围绕四条主线展开。第 \ref{sec:3-1} 节先建立 Banach 空间与 Hilbert 空间的几何工作台，说明完备性、正交分解和 Riesz 表示为何是后文分析与算法的基础接口。第 \ref{sec:3-2} 节把视角切换到弱拓扑与弱$^\ast$拓扑，给出 Banach--Alaoglu 定理以及自反空间中的弱紧性结论，从而为直接法与极小值存在准备真正可用的紧性语言。第 \ref{sec:3-3} 节再把矩阵转置升级为伴随算子，把约束映射、对偶变量与最优控制中的共轭系统统一到一个框架里。最后第 \ref{sec:3-4} 节只保留与优化直接相关的谱论工具：谱、紧算子、自伴算子、正算子以及紧自伴情形下的谱定理。

阅读本章时，需要始终区分三层环境。一般 Banach 空间足以承载凸泛函、对偶映射和有界线性算子；自反 Banach 空间开始提供弱紧性，这是存在性论证的关键台阶；Hilbert 空间则在此基础上额外拥有内积几何、正交投影与更强的谱工具，因此成为邻近算法、分裂方法和二次型分析的天然舞台。后文关于极小值存在、KKT 条件、Fenchel--Rockafellar 对偶与算子分裂算法的许多核心步骤，都会直接回链到本章。

\section{Banach 空间与 Hilbert 空间}

\label{sec:3-1}

无穷维优化首先面对的，不是某个具体算法，而是“变量究竟生活在哪个空间里”这一基础问题。有限维情形中，范数等价、闭有界集紧、线性泛函都能写成内积，这些事实使许多几何结构被自动隐藏。到了无穷维，赋范空间、Banach 空间与 Hilbert 空间之间的差别必须被明确区分：前者提供基本的线性与连续性框架，第二层加入完备性，第三层再加入内积几何与正交分解。后文中的邻近映射、投影算法、最小二乘和算子谱分解，都建立在这种层级区分之上。

\subsection{赋范空间与完备性}

\begin{definition}[赋范空间]\label{def:normed-space}
设 $X$ 为实或复线性空间。若映射 $\|\cdot\|\colon X\to [0,\infty)$ 满足：
\begin{enumerate}
  \item $\|x\|=0$ 当且仅当 $x=0$；
  \item 对任意标量 $\lambda$ 与任意 $x\in X$，有
  \[
  \|\lambda x\|=|\lambda|\,\|x\| ;
  \]
  \item 对任意 $x,y\in X$，有
  \[
  \|x+y\|\le \|x\|+\|y\| ,
  \]
\end{enumerate}
则称 $(X,\|\cdot\|)$ 为一个\emph{赋范空间}。由范数诱导的度量记为
\[
d(x,y):=\|x-y\|.
\]
\end{definition}

\begin{definition}[Banach 空间]\label{def:banach-space}
若赋范空间 $(X,\|\cdot\|)$ 在由定义~\ref{def:normed-space} 诱导的度量下完备，则称 $X$ 为一个\emph{Banach 空间}。
\end{definition}

完备性的作用，不只是保证 Cauchy 列有极限。更重要的是，它使闭图像、开映射、一致有界等基本定理可以成立，而这些结果又反过来控制算子方程、最优性条件与极值存在的稳定性。第 \ref{sec:1-2} 节中的 Baire 范畴方法，正是这种完备性力量的拓扑来源。

\begin{example}[函数空间的几何分层]\label{ex:lp-geometry}
对 $1\le p\le\infty$，空间 $\ell^p$ 与 $L^p$ 都是 Banach 空间；其中只有 $p=2$ 时，它们由内积
\[
\langle x,y\rangle=\sum_{n=1}^{\infty}x_n\overline{y_n},
\qquad
\langle f,g\rangle=\int f\overline{g}\,d\mu
\]
诱导范数，因此是 Hilbert 空间。把这个例子放在这里，是为了强调“完备”并不自动意味着“可做正交分解”；许多后文只在 Hilbert 空间成立的算法结论，在 $L^1$ 或 $\ell^\infty$ 中就必须改写。
\end{example}

\subsection{内积、正交与 Hilbert 空间}

\begin{definition}[Hilbert 空间]\label{def:hilbert-space}
设 $H$ 为实或复线性空间，$\langle\cdot,\cdot\rangle$ 为 $H$ 上的内积，范数由
\[
\|x\|:=\langle x,x\rangle^{1/2}
\]
给出。若 $(H,\|\cdot\|)$ 完备，则称 $H$ 为一个\emph{Hilbert 空间}。
\end{definition}

在 Hilbert 空间中，内积不仅给出范数，还把“方向”“垂直”“投影”“最短距离”这些欧氏概念全部带回了无穷维环境。设 $M\subset H$ 为非空子集，定义其正交补为
\[
M^\perp:=\{x\in H:\langle x,y\rangle=0,\ \forall y\in M\}.
\]
当 $M$ 是线性子空间时，$M^\perp$ 也是闭线性子空间，并且几何上对应于“所有与 $M$ 垂直的方向”。

\begin{proposition}[正交补的基本性质]
设 $H$ 为 Hilbert 空间，$M\subset H$ 为任意非空子集，则 $M^\perp$ 是 $H$ 的闭线性子空间。若 $M$ 本身是线性子空间，则
\[
\overline{M}\cap M^\perp=\{0\}.
\]
\end{proposition}

\begin{proof}
对任意 $y\in M$，映射 $x\mapsto \langle x,y\rangle$ 是连续线性泛函，因此集合
\[
\{x\in H:\langle x,y\rangle=0\}
\]
是闭线性子空间。于是
\[
M^\perp=\bigcap_{y\in M}\{x\in H:\langle x,y\rangle=0\}
\]
仍为闭线性子空间。若 $x\in \overline{M}\cap M^\perp$，取 $M$ 中列 $x_n\to x$，由内积连续性得
\[
\langle x,x\rangle=\lim_{n\to\infty}\langle x,x_n\rangle=0,
\]
故 $x=0$。
\end{proof}

\begin{theorem}[Hilbert 空间中的投影定理]\label{thm:orthogonal-projection}
设 $H$ 为 Hilbert 空间，$M\subset H$ 为非空闭线性子空间。则对任意 $x\in H$，存在唯一元素 $P_Mx\in M$ 使
\[
\|x-P_Mx\|=\inf_{y\in M}\|x-y\|.
\]
并且该极小元满足正交条件
\[
x-P_Mx\in M^\perp.
\]
因此每个 $x\in H$ 都可以唯一写成
\[
x=P_Mx+(x-P_Mx),\qquad P_Mx\in M,\ x-P_Mx\in M^\perp.
\]
\end{theorem}

\begin{proof}
记
\[
d:=\inf_{y\in M}\|x-y\|.
\]
取极小化列 $(y_n)\subset M$ 使 $\|x-y_n\|\to d$。由平行四边形恒等式，
\[
\left\|x-\frac{y_n+y_m}{2}\right\|^2+\left\|\frac{y_n-y_m}{2}\right\|^2
=\frac12\|x-y_n\|^2+\frac12\|x-y_m\|^2 .
\]
因为 $M$ 线性且闭，所以 $(y_n+y_m)/2\in M$，从而第一项不小于 $d^2$。于是
\[
\left\|\frac{y_n-y_m}{2}\right\|^2
\le
\frac12\|x-y_n\|^2+\frac12\|x-y_m\|^2-d^2\to 0.
\]
故 $(y_n)$ 是 Cauchy 列。由 $M$ 完备，存在 $y_\ast\in M$ 使 $y_n\to y_\ast$。连续性给出
\[
\|x-y_\ast\|=d,
\]
这证明了存在性。记 $P_Mx:=y_\ast$。

下证正交条件。任取 $z\in M$，对任意实数 $t$，因为 $P_Mx+tz\in M$，
\[
\|x-P_Mx\|^2\le \|x-P_Mx-tz\|^2
=\|x-P_Mx\|^2-2t\operatorname{Re}\langle x-P_Mx,z\rangle+t^2\|z\|^2.
\]
这对所有实数 $t$ 成立，只能说明
\[
\operatorname{Re}\langle x-P_Mx,z\rangle=0.
\]
在复 Hilbert 空间中再将 $z$ 替换为 $iz$，便得虚部也为零，因此
\[
\langle x-P_Mx,z\rangle=0,\qquad \forall z\in M.
\]
故 $x-P_Mx\in M^\perp$。

唯一性由 $\overline{M}\cap M^\perp=\{0\}$ 推出。若
\[
x=y_1+z_1=y_2+z_2,\qquad y_i\in M,\ z_i\in M^\perp,
\]
则 $y_1-y_2=z_2-z_1\in M\cap M^\perp$，从而 $y_1=y_2$、$z_1=z_2$。
\end{proof}

\paragraph{为何投影定理对算法重要}

\begin{remark}
定理~\ref{thm:orthogonal-projection} 是 Hilbert 空间算法几何的起点。后文的邻近算子、投影梯度法和算子分裂方法，本质上都依赖“最近点存在且可由正交条件刻画”这一结构。若只处在一般 Banach 空间，这种几何就不再自动存在，算法设计也会明显更微妙。
\end{remark}

\begin{example}[欧氏空间中的正交投影]\label{ex:euclidean-projection}
设 $H=\mathbb R^n$，$M$ 为某个线性子空间，取一组正交规范基组成矩阵 $Q$，则
\[
P_Mx=QQ^\top x.
\]
把这个例子放在这里，是为了把定理~\ref{thm:orthogonal-projection} 退回 Boyd 书中最熟悉的矩阵公式：无穷维投影定理并不是新对象，而是欧氏正交投影的真正上位版本。
\end{example}

\subsection{Hilbert 空间的 Riesz 表示定理}

\begin{theorem}[Hilbert 空间的 Riesz 表示定理]\label{thm:riesz-representation-hilbert}
设 $H$ 为 Hilbert 空间。则对每个连续线性泛函 $\varphi\in H^\ast$，存在唯一的元素 $y\in H$，使得
\[
\varphi(x)=\langle x,y\rangle,\qquad \forall x\in H.
\]
并且
\[
\|\varphi\|=\|y\|.
\]
因此映射
\[
J\colon H\to H^\ast,\qquad y\mapsto \langle\cdot,y\rangle
\]
是一个等距同构；在复 Hilbert 空间中，它是共轭线性的。
\end{theorem}

\begin{proof}
若 $\varphi=0$，取 $y=0$ 即可。现设 $\varphi\neq 0$。则 $\ker\varphi$ 是闭真子空间。由定理~\ref{thm:orthogonal-projection} 应用于 $M=\ker\varphi$，可得存在非零向量 $u\in (\ker\varphi)^\perp$。任取 $x\in H$，将其分解为
\[
x=z+\lambda u,\qquad z\in \ker\varphi,\ \lambda\in \mathbb F.
\]
因为 $z\in \ker\varphi$，有
\[
\varphi(x)=\lambda \varphi(u).
\]
另一方面，由 $u\perp z$ 得
\[
\lambda=\frac{\langle x,u\rangle}{\|u\|^2}.
\]
故
\[
\varphi(x)=\frac{\varphi(u)}{\|u\|^2}\langle x,u\rangle.
\]
若取
\[
y=
\begin{cases}
\dfrac{\overline{\varphi(u)}}{\|u\|^2}u, & \text{复 Hilbert 空间},\\[1ex]
\dfrac{\varphi(u)}{\|u\|^2}u, & \text{实 Hilbert 空间},
\end{cases}
\]
便有 $\varphi(x)=\langle x,y\rangle$ 对所有 $x$ 成立。

唯一性是显然的：若 $\langle x,y_1\rangle=\langle x,y_2\rangle$ 对所有 $x$ 成立，则
\[
\langle x,y_1-y_2\rangle=0,\qquad \forall x\in H.
\]
令 $x=y_1-y_2$ 即得 $y_1=y_2$。

最后证明范数恒等式。由 Cauchy--Schwarz 不等式，
\[
|\varphi(x)|=|\langle x,y\rangle|\le \|x\|\,\|y\|,
\]
故 $\|\varphi\|\le \|y\|$。反过来，若 $y\neq 0$，取 $x=y/\|y\|$，则
\[
\|\varphi\|\ge |\varphi(x)|=\left|\left\langle \frac{y}{\|y\|},y\right\rangle\right|=\|y\|.
\]
若 $y=0$ 则结论平凡。故 $\|\varphi\|=\|y\|$。
\end{proof}

\paragraph{Riesz 表示与后文伴随算子}

\begin{remark}
定理~\ref{thm:riesz-representation-hilbert} 的真正作用，是把 Hilbert 空间的对偶重新识别为原空间本身。第 \ref{sec:3-3} 节中的 Hilbert 伴随算子，正是通过这一识别把 Banach 空间中的对偶伴随 $T^\ast\colon Y^\ast\to X^\ast$ 拉回成 $T^\ast\colon H_2\to H_1$。
\end{remark}

\begin{example}[最小二乘与邻近映射触点]\label{ex:least-squares-hilbert}
在 Hilbert 空间 $H$ 中，给定闭线性子空间 $M$，最小二乘问题
\[
\min_{y\in M}\|x-y\|^2
\]
的解就是 $P_Mx$。更一般地，对闭凸集 $C$，最近点映射是后文邻近算子的几何原型。把这个例子放在这里，是为了说明第 9 章与第 10 章算法里的“prox”并不是凭空出现，而是投影定理在更一般凸泛函上的延伸。
\end{example}

\subsection*{有限维对照}

当 $H=\mathbb R^n$ 时，定义~\ref{def:hilbert-space} 与通常欧氏空间完全一致，定理~\ref{thm:orthogonal-projection} 退化为正交投影矩阵，定理~\ref{thm:riesz-representation-hilbert} 则退化为“每个线性泛函都可写成向量内积”。有限维凸优化中的最小二乘、法方程和正交分解，都只是本节结论在欧氏空间中的特例。

\subsection*{习题（待补）}

\begin{exercise}[投影定理的矩阵版占位]\label{exr:projection-matrix-placeholder}
补一题：把定理~\ref{thm:orthogonal-projection} 落回 $\mathbb R^n$，证明当 $M=\operatorname{range}(Q)$ 且 $Q$ 的列向量正交规范时，投影算子满足 $P_M=QQ^\top$。
\end{exercise}

\begin{exercise}[Banach 与 Hilbert 的桥接题占位]\label{exr:banach-hilbert-bridge-placeholder}
补一题：比较 $\ell^1$、$\ell^2$ 与 $\ell^\infty$ 中最近点问题的行为，说明为什么 Hilbert 空间中的正交几何不能无条件推广到一般 Banach 空间。
\end{exercise}


\section{弱拓扑与弱星拓扑}

\label{sec:3-2}

第 \ref{sec:3-1} 节建立了 Banach 与 Hilbert 的几何工作台，但仅靠范数拓扑并不足以处理无穷维极值存在。一个基本困难是：在无限维 Banach 空间中，闭有界集通常并不紧，因此“取极小化列，再抽收敛子列”的有限维套路会直接失效。要恢复紧性，我们必须引入比范数拓扑更弱的拓扑结构。弱拓扑与弱$^\ast$拓扑正是在这种背景下出现的：它们牺牲部分收敛强度，换来对偶球或有界集上的紧性，从而为直接法、对偶性和乘子理论提供可操作的极限机制。

\subsection{由函数族诱导的最弱拓扑}

设 $X$ 为集合，$\{f_i\}_{i\in I}$ 是一族映射 $f_i\colon X\to Y_i$，其中每个 $Y_i$ 都带有拓扑。使所有 $f_i$ 连续的最弱拓扑称为由这族函数诱导的\emph{初始拓扑}。弱拓扑与弱$^\ast$拓扑正是把所有连续线性泛函当作坐标函数所得到的初始拓扑。

\begin{definition}[弱拓扑]\label{def:weak-topology}
设 $X$ 为赋范空间。由对偶空间 $X^\ast$ 中全体连续线性泛函诱导的最弱拓扑，称为 $X$ 上的\emph{弱拓扑}，记为 $\sigma(X,X^\ast)$。换言之，弱拓扑是使每个
\[
\varphi\in X^\ast,\qquad \varphi\colon X\to \mathbb F
\]
都连续的最弱拓扑。
\end{definition}

\begin{definition}[弱星拓扑]\label{def:weak-star-topology}
设 $X$ 为赋范空间。由典范嵌入
\[
\hat{x}\colon X^\ast\to \mathbb F,\qquad \hat{x}(f)=f(x),\qquad x\in X,
\]
所诱导的 $X^\ast$ 上的最弱拓扑，称为 $X^\ast$ 上的\emph{弱$^\ast$拓扑}，记为 $\sigma(X^\ast,X)$。
\end{definition}

这两个拓扑的定义形式相近，但角色并不相同。弱拓扑使用的是 $X^\ast$ 中的全部连续线性泛函，因此是 $X$ 本身的拓扑；弱$^\ast$拓扑只使用来自原空间 $X$ 的评价泛函，因此是对偶空间 $X^\ast$ 上更弱的一种拓扑。Banach--Alaoglu 定理给出的紧性恰恰发生在后者上。

\paragraph{弱收敛为何必须用网}

\begin{remark}
由定义~\ref{def:weak-topology} 与定义~\ref{def:weak-star-topology} 可知，弱与弱$^\ast$收敛本质上是“所有坐标泛函逐点收敛”。但这些拓扑一般并非第一可数，因此不能安全地只谈序列。正如第 \ref{sec:1-1} 节定义~\ref{def:net-convergence} 所揭示的，一旦失去第一可数性，闭包和紧性就必须回到网语言。后文凡使用 Banach--Alaoglu 定理抽取极限点时，都应理解为“存在收敛子网”，而不是未经说明地退回子列。
\end{remark}

\begin{proposition}[弱收敛与弱星收敛的判别]
设 $X$ 为赋范空间。
\begin{enumerate}
  \item 网 $x_\alpha\to x$ 在 $\sigma(X,X^\ast)$ 下收敛，当且仅当对每个 $\varphi\in X^\ast$ 都有
  \[
  \varphi(x_\alpha)\to \varphi(x).
  \]
  \item 网 $f_\alpha\to f$ 在 $\sigma(X^\ast,X)$ 下收敛，当且仅当对每个 $x\in X$ 都有
  \[
  f_\alpha(x)\to f(x).
  \]
\end{enumerate}
\end{proposition}

\begin{proof}
这是初始拓扑的直接刻画：网在初始拓扑下收敛，当且仅当其经所有坐标映射后的像网都收敛。这里的坐标映射分别是 $\varphi\in X^\ast$ 与 $\hat{x}\in (X^\ast)^\ast$。
\end{proof}

\subsection{典范嵌入与自反性}

\begin{definition}[自反空间]\label{def:reflexive-space}
设 $X$ 为赋范空间。定义典范嵌入
\[
J_X\colon X\to X^{\ast\ast},\qquad J_X(x)(\varphi)=\varphi(x),\qquad \varphi\in X^\ast.
\]
若 $J_X$ 为满射，则称 $X$ 为\emph{自反空间}。
\end{definition}

\begin{proposition}[典范嵌入是等距嵌入]
设 $X$ 为赋范空间，则 $J_X$ 线性且满足
\[
\|J_X(x)\|=\|x\|,\qquad \forall x\in X.
\]
因此 $X$ 可以等距地看成 $X^{\ast\ast}$ 的闭子空间。
\end{proposition}

\begin{proof}
线性由定义直接得到。对任意 $x\in X$，
\[
\|J_X(x)\|
=\sup_{\|\varphi\|\le 1}|J_X(x)(\varphi)|
=\sup_{\|\varphi\|\le 1}|\varphi(x)|
\le \|x\|.
\]
反向不等式可由 Hahn--Banach 定理得到：对任意 $x\neq 0$，存在 $\varphi\in X^\ast$ 使 $\|\varphi\|=1$ 且 $\varphi(x)=\|x\|$，于是 $\|J_X(x)\|\ge \|x\|$。故二者相等。
\end{proof}

\paragraph{自反性的优化含义}

\begin{remark}
自反性不是纯粹的空间分类术语。它意味着原空间的有界集能通过典范嵌入进入双对偶的弱$^\ast$紧框架，再拉回到原空间自身，从而把 Banach--Alaoglu 定理转化为原空间中的弱紧性。这正是无穷维极小值存在论证里最常用的路径。
\end{remark}

\subsection{Banach--Alaoglu 定理与弱紧性}

前面的讨论已经把弱$^\ast$拓扑放到了正确位置：它是保留对偶球紧性的那层拓扑。下面给出 Banach--Alaoglu 定理及其标准证明。证明的核心思路只有两步：先把对偶球嵌入一个逐点坐标的乘积紧空间，再证明其像集在该乘积空间中闭。

\begin{theorem}[Banach--Alaoglu 定理]\label{thm:banach-alaoglu}
设 $X$ 为赋范空间，则对偶空间单位闭球
\[
B_{X^\ast}:=\{f\in X^\ast:\|f\|\le 1\}
\]
在弱$^\ast$拓扑 $\sigma(X^\ast,X)$ 下是紧的。
\end{theorem}

\begin{proof}
对每个 $x\in X$，记
\[
D_x:=\{z\in \mathbb F:|z|\le \|x\|\}.
\]
考虑映射
\[
\Phi\colon B_{X^\ast}\to \prod_{x\in X}D_x,\qquad \Phi(f)=(f(x))_{x\in X}.
\]
积空间按 Tychonoff 定理是紧的，因为每个 $D_x$ 都是紧集。弱$^\ast$拓扑正是使所有坐标映射
\[
f\mapsto f(x)
\]
连续的最弱拓扑，因此 $\Phi$ 是一个拓扑嵌入。只需证明 $\Phi(B_{X^\ast})$ 在积空间中是闭集。

设 $\Phi(f_\alpha)\to (z_x)_{x\in X}$ 于积空间。由坐标收敛知对每个 $x\in X$，$f_\alpha(x)\to z_x$。定义
\[
f(x):=z_x.
\]
极限保线性，因此 $f$ 线性。再由
\[
|f(x)|=\lim_\alpha |f_\alpha(x)|\le \|x\|
\]
可知 $f$ 有界且 $\|f\|\le 1$。故 $f\in B_{X^\ast}$ 且 $\Phi(f)=(z_x)_{x\in X}$。这说明像集闭，因而是紧的，从而 $B_{X^\ast}$ 在弱$^\ast$拓扑下紧。
\end{proof}

\begin{proposition}[自反空间中的闭球弱紧]
设 $X$ 为自反 Banach 空间，则闭单位球
\[
B_X:=\{x\in X:\|x\|\le 1\}
\]
在弱拓扑 $\sigma(X,X^\ast)$ 下紧。
\end{proposition}

\begin{proof}
由定义~\ref{def:reflexive-space}，典范嵌入 $J_X$ 把 $X$ 等距同构到 $X^{\ast\ast}$。根据定理~\ref{thm:banach-alaoglu}，$B_{X^{\ast\ast}}$ 在 $\sigma(X^{\ast\ast},X^\ast)$ 下弱$^\ast$紧。由于 $J_X$ 在自反情形下满射，它把 $B_X$ 与 $B_{X^{\ast\ast}}$ 对应起来，而 $\sigma(X,X^\ast)$ 恰是从双对偶弱$^\ast$拓扑拉回得到的拓扑，故 $B_X$ 弱紧。
\end{proof}

\begin{corollary}[Hilbert 空间中的弱紧性]\label{cor:hilbert-weak-compactness}
设 $H$ 为 Hilbert 空间，则每个非空闭有界凸集在弱拓扑下都是紧的；特别地，每个闭球
\[
\{x\in H:\|x\|\le r\}
\]
在弱拓扑下紧。
\end{corollary}

\begin{proof}
第 \ref{sec:3-1} 节定理~\ref{thm:riesz-representation-hilbert} 说明 $H^\ast$ 与 $H$ 等距同构，因此 $H$ 自反。由前一命题，闭球弱紧。任意闭有界凸集 $C\subset H$ 包含于某个闭球中；由于 $C$ 在范数拓扑下闭，而弱拓扑比范数拓扑更弱，范数闭凸集在 Hilbert 空间中也是弱闭集。紧空间的闭子集仍紧，故 $C$ 弱紧。
\end{proof}

\paragraph{极小值存在的正确表述}

\begin{remark}
推论~\ref{cor:hilbert-weak-compactness} 并不意味着“只要问题在 Hilbert 空间里就一定有极小值”。正确的说法是：若可行集在弱拓扑下紧，且目标泛函对弱拓扑下半连续，那么直接法才可能成立。实际应用中，弱紧性通常来自三种来源：有界闭凸可行集、自反性、以及由强制性把极小化序列压进某个弱紧闭球。缺少这些条件时，极小值完全可能不存在。
\end{remark}

\begin{example}[最小抽象例子：对偶球的弱星紧性]\label{ex:dual-ball-weak-star}
设 $X=C([0,1])$，则 $X^\ast$ 可以通过第 \ref{sec:2-3} 节定理~\ref{thm:riesz-representation} 识别为有限符号 Radon 测度空间。定理~\ref{thm:banach-alaoglu} 表明其中单位球在弱$^\ast$拓扑下紧。把这个例子放在这里，是为了说明对偶变量并不一定是向量；在无穷维中，它们往往是测度，而其紧性首先来自弱$^\ast$拓扑。
\end{example}

\begin{example}[有限维坍缩：闭有界即紧]\label{ex:finite-dimensional-compactness}
若 $X=\mathbb R^n$，则弱拓扑、弱$^\ast$拓扑与通常欧氏拓扑一致，因此定理~\ref{thm:banach-alaoglu} 与推论~\ref{cor:hilbert-weak-compactness} 都退化为 Heine--Borel 定理。把这个例子放在这里，是为了说明 Boyd 书里大量“不言自明”的存在性步骤，实际上依赖了维数有限这一隐藏条件。
\end{example}

\begin{example}[算法触点：弱极限与直接法]\label{ex:weak-direct-method}
在 Hilbert 空间上考虑泛函
\[
u\mapsto \frac12\|Au-b\|^2+\lambda \|u\|^2,
\]
若 $A$ 为有界线性算子，则该泛函在闭球上的极小化常先通过推论~\ref{cor:hilbert-weak-compactness} 取到弱收敛子网，再结合弱下半连续性得到极小元。把这个例子放在这里，是为了给后文算法收敛与存在性证明提供统一模板：先有界，再弱紧，再过极限。
\end{example}

\subsection*{有限维对照}

在有限维空间中，弱拓扑与通常拓扑没有区别，弱$^\ast$拓扑也与对偶空间上的通常拓扑一致。因此很多有限维教材会把“闭有界”“紧”“对偶球紧”混在一起说，而不会显式引入定义~\ref{def:weak-topology} 与定义~\ref{def:weak-star-topology}。本节的作用，是在无穷维情形下把这些层次重新分开，并用定理~\ref{thm:banach-alaoglu} 指出真正保住紧性的拓扑位置。

\subsection*{习题（待补）}

\begin{exercise}[弱拓扑与网占位]\label{exr:weak-topology-net-placeholder}
补一题：结合第 \ref{sec:1-1} 节命题~\ref{prop:closure-net}，说明为什么 Banach--Alaoglu 定理的结论自然应表述为“任意网有收敛子网”，而不是“任意序列有收敛子列”。
\end{exercise}

\begin{exercise}[直接法桥接题占位]\label{exr:direct-method-placeholder}
补一题：在 Hilbert 空间中证明闭球上任意连续凸泛函若再具弱下半连续性，则存在极小元，并指出该证明中哪些步骤依赖推论~\ref{cor:hilbert-weak-compactness}。
\end{exercise}


\section{线性算子与伴随算子}

\label{sec:3-3}

从优化的角度看，矩阵只是线性算子在有限维空间中的坐标表达。进入无穷维以后，线性约束 $Ax=b$、梯度映射、积分方程、偏微分方程和状态转移算子都必须用有界线性算子的语言统一处理。尤其重要的是：有限维里熟悉的转置 $A^\top$，在一般 Banach 空间中不再直接作用在原空间上，而是首先作用在对偶空间上，这就是伴随算子的真正来源。只有把这一步说清，后文的 Fenchel 对偶、KKT 乘子和最优控制中的伴随方程才有严密的对象归属。

\subsection{有界线性算子与算子空间}

\begin{definition}[有界线性算子]\label{def:bounded-linear-operator}
设 $X,Y$ 为赋范空间。线性映射 $T\colon X\to Y$ 若存在常数 $C\ge 0$ 使
\[
\|Tx\|_Y\le C\|x\|_X,\qquad \forall x\in X,
\]
则称 $T$ 为一个\emph{有界线性算子}。全体此类算子记为 $\mathcal B(X,Y)$，其范数定义为
\[
\|T\|:=\sup_{\|x\|_X\le 1}\|Tx\|_Y.
\]
\end{definition}

在赋范空间之间，线性算子的连续性与有界性是等价的，因此定义~\ref{def:bounded-linear-operator} 同时给出了“连续线性算子”的全部内容。对优化而言，这意味着线性约束、梯度线性化、状态方程和积分变换都可以统一地放进同一个 Banach 空间 $\mathcal B(X,Y)$ 中讨论。

\begin{proposition}[算子范数的基本性质]
设 $X,Y,Z$ 为赋范空间，$S,T\in \mathcal B(X,Y)$，$R\in \mathcal B(Y,Z)$，则
\begin{enumerate}
  \item $\mathcal B(X,Y)$ 是线性空间；
  \item 对任意 $x\in X$，有
  \[
  \|Tx\|_Y\le \|T\|\,\|x\|_X ;
  \]
  \item 复合满足
  \[
  \|RT\|\le \|R\|\,\|T\|.
  \]
\end{enumerate}
\end{proposition}

\begin{proof}
前两条直接来自定义~\ref{def:bounded-linear-operator}。对第三条，任取 $\|x\|_X\le 1$，则
\[
\|RTx\|_Z\le \|R\|\,\|Tx\|_Y\le \|R\|\,\|T\|\,\|x\|_X\le \|R\|\,\|T\|.
\]
对单位球取上确界即得。
\end{proof}

\begin{proposition}[Banach 空间值算子空间的完备性]\label{prop:operator-space-banach}
设 $X$ 为赋范空间，$Y$ 为 Banach 空间，则 $\mathcal B(X,Y)$ 配上算子范数后也是 Banach 空间。
\end{proposition}

\begin{proof}
设 $(T_n)$ 是 $\mathcal B(X,Y)$ 中的 Cauchy 列。对任意固定 $x\in X$，
\[
\|T_nx-T_mx\|_Y\le \|T_n-T_m\|\,\|x\|_X,
\]
故 $(T_nx)$ 在 $Y$ 中是 Cauchy 列。由于 $Y$ 完备，存在元素 $Tx\in Y$ 使 $T_nx\to Tx$。于是定义了映射 $T\colon X\to Y$。

由极限与线性运算交换，$T$ 线性。再由
\[
\|Tx\|_Y=\lim_{n\to\infty}\|T_nx\|_Y
\le \left(\sup_n \|T_n\|\right)\|x\|_X
\]
可知 $T$ 有界。最后，因为 $(T_n)$ 在算子范数下是 Cauchy 列，给定 $\varepsilon>0$，取 $N$ 使 $n,m\ge N$ 时 $\|T_n-T_m\|<\varepsilon$。固定 $n\ge N$ 并令 $m\to\infty$，得到
\[
\|T_n-T\|\le \varepsilon.
\]
故 $T_n\to T$ 于算子范数，证明完毕。
\end{proof}

\subsection{Banach 空间中的对偶伴随}

\begin{definition}[伴随算子]\label{def:adjoint-operator}
设 $T\in \mathcal B(X,Y)$，其中 $X,Y$ 为赋范空间。定义映射
\[
T^\ast\colon Y^\ast\to X^\ast,\qquad T^\ast f:=f\circ T.
\]
称 $T^\ast$ 为 $T$ 的\emph{对偶伴随算子}。
\end{definition}

\begin{proposition}[伴随算子的基本性质]\label{prop:adjoint-basic-properties}
设 $X,Y,Z$ 为赋范空间，$T,S\in \mathcal B(X,Y)$，$R\in \mathcal B(Y,Z)$，则：
\begin{enumerate}
  \item $T^\ast\in \mathcal B(Y^\ast,X^\ast)$ 且
  \[
  \|T^\ast\|=\|T\| ;
  \]
  \item $(T+S)^\ast=T^\ast+S^\ast$；
  \item $(\lambda T)^\ast=\lambda T^\ast$；
  \item $(RT)^\ast=T^\ast R^\ast$。
\end{enumerate}
\end{proposition}

\begin{proof}
对任意 $f\in Y^\ast$，
\[
\|T^\ast f\|
=\sup_{\|x\|_X\le 1}|f(Tx)|
\le \|f\|\,\|T\|,
\]
故 $T^\ast$ 有界且 $\|T^\ast\|\le \|T\|$。反过来，任取 $\varepsilon>0$，可取 $\|x\|_X=1$ 使
\[
\|Tx\|_Y\ge \|T\|-\varepsilon.
\]
由 Hahn--Banach 定理，存在 $\|f\|=1$ 使
\[
|f(Tx)|=\|Tx\|_Y.
\]
于是
\[
\|T^\ast\|
\ge \|T^\ast f\|
\ge |(T^\ast f)(x)|
=|f(Tx)|
\ge \|T\|-\varepsilon.
\]
令 $\varepsilon\downarrow 0$ 得 $\|T^\ast\|=\|T\|$。其余三条都由定义~\ref{def:adjoint-operator} 直接验证，例如
\[
(RT)^\ast g=g\circ R\circ T=T^\ast(R^\ast g).
\]
\end{proof}

\paragraph{核、像与最优性条件}

\begin{remark}
当约束写成 $Tu=b$ 时，拉格朗日函数中的乘子自然属于 $Y^\ast$，而驻点条件则涉及 $T^\ast\lambda\in X^\ast$。因此在 Banach 空间里，伴随首先是“把对偶变量拉回原问题空间的对偶空间”的算子，而不是像有限维那样仍留在同一个坐标空间中。
\end{remark}

\subsection{Hilbert 空间中的伴随}

若 $X,Y$ 都是 Hilbert 空间，则第 \ref{sec:3-1} 节定理~\ref{thm:riesz-representation-hilbert} 可以把 $X^\ast$ 与 $X$、$Y^\ast$ 与 $Y$ 重新识别起来，于是 Banach 空间中的对偶伴随可被改写成作用在原 Hilbert 空间之间的伴随。

\begin{corollary}[Hilbert 空间中的伴随表示]\label{cor:hilbert-adjoint-via-riesz}
设 $H_1,H_2$ 为 Hilbert 空间，$T\in \mathcal B(H_1,H_2)$。则存在唯一算子
\[
T^\ast\in \mathcal B(H_2,H_1)
\]
使得
\[
\langle Tx,y\rangle_{H_2}=\langle x,T^\ast y\rangle_{H_1},
\qquad \forall x\in H_1,\ \forall y\in H_2.
\]
并且
\[
\|T^\ast\|=\|T\|.
\]
\end{corollary}

\begin{proof}
固定 $y\in H_2$，考虑映射
\[
\varphi_y\colon H_1\to \mathbb F,\qquad \varphi_y(x):=\langle Tx,y\rangle_{H_2}.
\]
它是 $H_1$ 上的连续线性泛函，并满足
\[
|\varphi_y(x)|\le \|T\|\,\|y\|\,\|x\|.
\]
由定理~\ref{thm:riesz-representation-hilbert}，存在唯一元素 $z\in H_1$ 使
\[
\varphi_y(x)=\langle x,z\rangle_{H_1},\qquad \forall x\in H_1.
\]
定义 $T^\ast y:=z$，就得到所需等式。线性与有界性可由 Riesz 表示中的唯一性直接推出，且
\[
\|T^\ast\|=\|T\|
\]
来自命题~\ref{prop:adjoint-basic-properties} 与 Riesz 同构的等距性。
\end{proof}

\begin{example}[矩阵转置作为伴随]\label{ex:matrix-transpose-adjoint}
当 $H_1=\mathbb R^m$、$H_2=\mathbb R^n$ 配标准内积时，推论~\ref{cor:hilbert-adjoint-via-riesz} 中的 $T^\ast$ 正是矩阵转置 $A^\top$。把这个例子放在这里，是为了说明本节不是引入一种“比矩阵更抽象的新记号”，而是在解释矩阵转置在无穷维环境中的真实含义。
\end{example}

\begin{example}[积分算子的伴随]\label{ex:integral-operator-adjoint}
设
\[
(Tf)(s)=\int_0^1 K(s,t)f(t)\,dt,
\qquad f\in L^2(0,1),
\]
其中核函数 $K\in L^2((0,1)^2)$。则 $T$ 是 $L^2(0,1)$ 上的有界线性算子，并满足
\[
(T^\ast g)(t)=\int_0^1 \overline{K(s,t)}\,g(s)\,ds.
\]
把这个例子放在这里，是为了展示伴随算子在函数空间中仍然是可计算对象；后文对偶系统与 Euler--Lagrange 型条件里经常出现的正是这种核的反向传播。
\end{example}

\begin{example}[控制与 PDE 中的伴随方程]\label{ex:pde-adjoint}
设状态方程形式上写成
\[
\mathcal A u = f
\]
并以某个 Hilbert 空间作为状态空间。当目标泛函对状态求导后，需要把变化量从状态端拉回控制端时，出现的正是 $\mathcal A^\ast$。例如在抛物型控制问题中，前向方程与反向伴随方程组成一对耦合系统。把这个例子放在这里，是为了说明“伴随”并不只是线性代数术语，而是最优控制中梯度计算的核心接口。
\end{example}

\subsection*{有限维对照}

在 Boyd 书里，线性约束与 KKT 条件通常写成 $A^\top \lambda$。本节说明，这个符号在无穷维中的真实上位版本是定义~\ref{def:adjoint-operator} 中的 $T^\ast f=f\circ T$；只有在 Hilbert 空间并借助推论~\ref{cor:hilbert-adjoint-via-riesz} 时，它才重新坍缩为熟悉的“转置仍作用在原空间上”。

\subsection*{习题（待补）}

\begin{exercise}[积分核伴随占位]\label{exr:integral-adjoint-placeholder}
补一题：对例~\ref{ex:integral-operator-adjoint} 中的积分算子进行严格计算，证明给出的 $T^\ast$ 确实满足推论~\ref{cor:hilbert-adjoint-via-riesz} 的定义恒等式。
\end{exercise}

\begin{exercise}[有限维桥接题占位]\label{exr:transpose-bridge-placeholder}
补一题：把命题~\ref{prop:adjoint-basic-properties} 退回矩阵情形，逐条验证它们恰好对应于 $(A+B)^\top=A^\top+B^\top$、$(\lambda A)^\top=\lambda A^\top$ 与 $(BA)^\top=A^\top B^\top$。
\end{exercise}


\section{算子代数初步}

\label{sec:3-4}

本节的标题虽然包含“算子代数”，但目标并不是进入一般 $C^\ast$ 代数理论，而是提炼出优化真正会反复调用的谱论工具。对二次泛函、Hessian 型算子、协方差算子、核积分算子和半正定约束来说，最关键的问题只有几个：一个有界算子的谱是什么，紧算子为何像“无穷维的近似有限秩矩阵”，自伴与正性如何把几何信息编码进二次型，以及在什么条件下算子可以被“对角化”。这些问题在紧自伴情形下有最清晰的答案，而这正是本节要建立的工作台。

\subsection{谱与谱半径}

\begin{definition}[有界算子的谱]\label{def:spectrum-bounded-operator}
设 $X$ 为复 Banach 空间，$T\in \mathcal B(X)$。若复数 $\lambda$ 使算子
\[
\lambda I-T
\]
不可逆，则称 $\lambda$ 属于 $T$ 的\emph{谱}，记为 $\sigma(T)$。其补集
\[
\rho(T):=\mathbb C\setminus \sigma(T)
\]
称为 $T$ 的\emph{预解集}。谱半径定义为
\[
r(T):=\sup\{|\lambda|:\lambda\in \sigma(T)\}.
\]
\end{definition}

\begin{proposition}[谱的基本性质]
设 $X$ 为复 Banach 空间，$T\in \mathcal B(X)$。则 $\sigma(T)$ 是非空紧集，且
\[
r(T)\le \|T\|.
\]
\end{proposition}

\begin{proof}
当 $|\lambda|>\|T\|$ 时，
\[
\lambda I-T=\lambda\left(I-\lambda^{-1}T\right),
\]
且 $\|\lambda^{-1}T\|<1$。由 Neumann 级数
\[
\left(I-\lambda^{-1}T\right)^{-1}
=\sum_{n=0}^{\infty}\lambda^{-n}T^n
\]
可知 $\lambda I-T$ 可逆，所以
\[
\sigma(T)\subset \{\lambda\in \mathbb C:|\lambda|\le \|T\|\}.
\]
另一方面，可逆元在 Banach 代数中构成开集，因此预解集 $\rho(T)$ 开，谱 $\sigma(T)$ 闭，于是谱为有界闭集，从而紧。

至于非空性，这是 Banach 代数谱论的基本事实：若 $\sigma(T)=\varnothing$，则预解算子
\[
\lambda\mapsto (\lambda I-T)^{-1}
\]
成为整函数，并满足无穷远处的衰减估计，借助 Liouville 定理会推出它恒为零，矛盾。故谱非空。
\end{proof}

\paragraph{为何谱语言对优化有用}

\begin{remark}
对有限维矩阵，谱不过是特征值集合；但在无穷维中，谱提供了比特征值更稳定的语言。后文讨论二次型、稳定性与半正定约束时，真正不变的对象是谱，而不是某个固定坐标下的矩阵表示。
\end{remark}

\subsection{紧算子、自伴算子与正算子}

\begin{definition}[紧算子]\label{def:compact-operator}
设 $X,Y$ 为 Banach 空间。若 $T\in \mathcal B(X,Y)$ 满足：$X$ 的任意有界集在 $T$ 作用下的像都是 $Y$ 中相对紧集，则称 $T$ 为一个\emph{紧算子}。
\end{definition}

\begin{definition}[自伴算子]\label{def:self-adjoint-operator}
设 $H$ 为 Hilbert 空间。基于第 \ref{sec:3-3} 节定义~\ref{def:adjoint-operator} 以及推论~\ref{cor:hilbert-adjoint-via-riesz}，若算子 $T\in \mathcal B(H)$ 满足
\[
T^\ast=T,
\]
则称 $T$ 为\emph{自伴算子}。
\end{definition}

\begin{definition}[正算子]\label{def:positive-operator}
设 $H$ 为 Hilbert 空间。若 $T\in \mathcal B(H)$ 满足
\[
\langle Tx,x\rangle\ge 0,\qquad \forall x\in H,
\]
则称 $T$ 为\emph{正算子}，记作 $T\ge 0$。
\end{definition}

\begin{proposition}[紧算子把弱收敛提升为强收敛]
设 $X,Y$ 为 Banach 空间，$T\in \mathcal B(X,Y)$ 为紧算子。若有界网 $x_\alpha$ 在 $X$ 中弱收敛到 $x$，则
\[
Tx_\alpha\to Tx
\]
在 $Y$ 的范数下成立。
\end{proposition}

\begin{proof}
因为 $x_\alpha\rightharpoonup x$，网 $(x_\alpha)$ 有界。由定义~\ref{def:compact-operator}，像网 $(Tx_\alpha)$ 落在某个紧集的闭包中，因此任意子网都有范数收敛子网。设 $Tx_\beta\to y$ 于范数，则它也弱收敛到 $y$。另一方面，$T$ 连续且线性，故 $Tx_\beta\rightharpoonup Tx$。弱极限唯一，遂 $y=Tx$。因此 $(Tx_\alpha)$ 的每个收敛子网都只能收敛到 $Tx$。紧 Hausdorff 空间中的网若所有收敛子网极限相同，则原网收敛到该极限，故结论成立。
\end{proof}

\begin{proposition}[正算子的基本性质]
设 $H$ 为 Hilbert 空间，$T\in \mathcal B(H)$。
\begin{enumerate}
  \item 若 $T\ge 0$，则 $T$ 必为自伴算子；
  \item 若 $T$ 自伴，则 $\langle Tx,x\rangle\in \mathbb R$ 对任意 $x\in H$ 成立；
  \item 若 $T$ 自伴且 $\lambda\in \sigma(T)$ 为特征值，则 $\lambda\in \mathbb R$。
\end{enumerate}
\end{proposition}

\begin{proof}
第二条直接来自
\[
\langle Tx,x\rangle=\langle x,Tx\rangle=\overline{\langle Tx,x\rangle}.
\]
第三条则由 $Tx=\lambda x$ 推出
\[
\lambda \|x\|^2=\langle Tx,x\rangle\in \mathbb R.
\]

对第一条，考虑极化恒等式。实 Hilbert 空间中有
\[
4\langle Tx,y\rangle
=\langle T(x+y),x+y\rangle-\langle T(x-y),x-y\rangle.
\]
复 Hilbert 空间中有
\[
4\langle Tx,y\rangle
=\sum_{k=0}^{3} i^k\langle T(x+i^ky),x+i^ky\rangle.
\]
由于 $T\ge 0$ 时右端均可由二次型
\[
q(z):=\langle Tz,z\rangle
\]
完全恢复，得到
\[
\langle Tx,y\rangle=\overline{\langle Ty,x\rangle}=\langle x,Ty\rangle,
\]
故 $T=T^\ast$。
\end{proof}

\begin{proposition}[正算子的范数刻画]
设 $H$ 为 Hilbert 空间，$T\in \mathcal B(H)$ 且 $T\ge 0$。则
\[
\|T\|=\sup_{\|x\|=1}\langle Tx,x\rangle.
\]
\end{proposition}

\begin{proof}
对任意单位向量 $x$，显然有
\[
\langle Tx,x\rangle\le \|Tx\|\,\|x\|\le \|T\|,
\]
故右端上确界不超过 $\|T\|$。

反过来，对正算子定义半内积
\[
[u,v]_T:=\langle Tu,v\rangle.
\]
由 Cauchy--Schwarz 不等式，
\[
|[x,Tx]_T|^2\le [x,x]_T\,[Tx,Tx]_T.
\]
亦即
\[
\|Tx\|^4
\le \langle Tx,x\rangle \,\langle T^2x,Tx\rangle
\le \langle Tx,x\rangle\,\|T^2x\|\,\|Tx\|
\le \|T\|\,\langle Tx,x\rangle\,\|Tx\|^2.
\]
若 $Tx\neq 0$，可约去 $\|Tx\|^2$，得到
\[
\|Tx\|^2\le \|T\|\,\langle Tx,x\rangle.
\]
对单位向量取上确界，记
\[
a:=\sup_{\|x\|=1}\langle Tx,x\rangle,
\]
便有
\[
\|T\|^2=\sup_{\|x\|=1}\|Tx\|^2\le \|T\|\,a.
\]
若 $T\neq 0$，则 $\|T\|>0$，故 $\|T\|\le a$；若 $T=0$ 则结论显然。综上 $\|T\|=a$。
\end{proof}

\begin{example}[有限维坍缩：对称矩阵与半正定矩阵]\label{ex:symmetric-matrix-spectrum}
当 $H=\mathbb R^n$ 时，定义~\ref{def:self-adjoint-operator} 退化为对称矩阵，定义~\ref{def:positive-operator} 退化为半正定矩阵。把这个例子放在这里，是为了说明半正定规划中的矩阵锥并不是一个孤立对象，而是 Hilbert 空间正算子锥的有限维切片。
\end{example}

\begin{example}[紧算子的典型原型：积分算子]\label{ex:compact-integral-operator}
在 $L^2(0,1)$ 上，若核函数 $K$ 连续，则
\[
(Tf)(s)=\int_0^1 K(s,t)f(t)\,dt
\]
定义了一个紧算子。若再满足 $K(s,t)=\overline{K(t,s)}$，则 $T$ 还是自伴算子。把这个例子放在这里，是为了把本节的抽象谱理论直接连接到后文的协方差算子、核方法和 Fredholm 型方程。
\end{example}

\subsection{紧自伴算子的谱定理}

\begin{lemma}[正紧自伴算子的主特征值]
设 $H$ 为 Hilbert 空间，$T\in \mathcal B(H)$ 为非零、正、紧、自伴算子。则存在单位向量 $e\in H$ 与实数 $\mu>0$，使
\[
Te=\mu e,
\qquad
\mu=\|T\|.
\]
\end{lemma}

\begin{proof}
考虑单位闭球
\[
B:=\{x\in H:\|x\|\le 1\}.
\]
由推论~\ref{cor:hilbert-weak-compactness}，$B$ 在弱拓扑下紧。又由紧算子把弱收敛提升为强收敛这一命题，函数
\[
q(x):=\langle Tx,x\rangle
\]
在 $B$ 上是弱连续的：若 $x_\alpha\rightharpoonup x$，则弱收敛网必有界，且 $Tx_\alpha\to Tx$ 于范数，于是
\[
\langle Tx_\alpha,x_\alpha\rangle-\langle Tx,x\rangle
=\langle Tx_\alpha-Tx,x_\alpha\rangle+\langle Tx,x_\alpha-x\rangle\to 0.
\]
因此 $q$ 在 $B$ 上取到最大值。取单位向量 $e$ 满足
\[
q(e)=\max_{\|x\|\le 1} q(x).
\]
由于 $T\ge 0$ 且 $T\neq 0$，有 $q(e)>0$。记 $\mu:=q(e)$。

对任意 $y\in H$ 且 $\langle y,e\rangle=0$，考虑函数
\[
\phi(t):=q\!\left(\frac{e+ty}{\|e+ty\|}\right),\qquad t\in \mathbb R.
\]
因为 $e$ 在单位球上使 $q$ 取极大值，故 $\phi'(0)=0$。直接计算得到
\[
\operatorname{Re}\langle Te,y\rangle=0.
\]
在复 Hilbert 空间中再将 $y$ 换为 $iy$，可得虚部也为零，因此
\[
\langle Te,y\rangle=0,\qquad \forall y\perp e.
\]
这说明 $Te$ 与所有垂直于 $e$ 的向量都正交，故 $Te$ 必与 $e$ 共线，即存在标量 $\lambda$ 使 $Te=\lambda e$。再与 $e$ 取内积，得到
\[
\lambda=\langle Te,e\rangle=\mu.
\]
最后由上一命题，
\[
\|T\|
=\sup_{\|x\|\le 1}\langle Tx,x\rangle
=q(e)=\mu
\]
得出 $\mu=\|T\|$。
\end{proof}

\begin{proposition}[紧算子的非零特征值只能有限重]
设 $H$ 为 Hilbert 空间，$T\in \mathcal B(H)$ 为紧算子。若 $\lambda\neq 0$ 是 $T$ 的特征值，则对应特征子空间
\[
\ker(T-\lambda I)
\]
必为有限维。
\end{proposition}

\begin{proof}
若该特征子空间无限维，则可以从中取一组标准正交列 $(e_n)$。对每个 $n$，
\[
Te_n=\lambda e_n.
\]
由于 $\lambda\neq 0$，序列 $(Te_n)$ 的任意两项之间距离都等于
\[
\|Te_n-Te_m\|=|\lambda|\,\|e_n-e_m\|=\sqrt{2}\,|\lambda|,\qquad n\neq m,
\]
故不可能含有收敛子列。这与 $T$ 为紧算子矛盾。
\end{proof}

\begin{theorem}[紧自伴算子的谱定理]\label{thm:compact-selfadjoint-spectral}
设 $H$ 为 Hilbert 空间，$T\in \mathcal B(H)$ 为紧自伴算子。则存在至多可数的标准正交特征向量族 $\{e_n\}$ 与实特征值列 $\{\lambda_n\}$，满足：
\begin{enumerate}
  \item 对每个 $n$，
  \[
  Te_n=\lambda_n e_n;
  \]
  \item 若非零特征值无限多个，则
  \[
  \lambda_n\to 0;
  \]
  \item 对任意 $x\in H$，
  \[
  Tx=\sum_{n} \lambda_n \langle x,e_n\rangle e_n,
  \]
  其中级数在 $H$ 中按范数收敛；
  \item 非零谱恰好由这些非零特征值组成，每个非零特征值重数有限，且 $0$ 是唯一可能的聚点。
\end{enumerate}
\end{theorem}

\begin{proof}
若 $T=0$，取空特征族即可。以下设 $T\neq 0$。

先考虑
\[
S:=T^2=T^\ast T.
\]
则 $S$ 是正、紧、自伴算子。由上一个引理，存在单位向量 $u$ 与标量 $\mu=\|S\|=\|T\|^2>0$，使
\[
Su=\mu u.
\]
记
\[
E_\mu:=\ker(S-\mu I).
\]
由上一命题，$E_\mu$ 有限维。又因为 $TS=ST$，所以 $T(E_\mu)\subseteq E_\mu$。将 $T$ 限制到有限维 Hilbert 空间 $E_\mu$ 上，便得到一个有限维自伴算子；由有限维线性代数，它可以正交对角化。因此存在单位向量 $e_1\in E_\mu$ 与实数 $\lambda_1$，使
\[
Te_1=\lambda_1 e_1,
\qquad
|\lambda_1|=\|T\|.
\]

接着做递归构造。已取到两两正交的单位特征向量 $e_1,\dots,e_n$ 后，记
\[
M_n:=\operatorname{span}\{e_1,\dots,e_n\}^\perp.
\]
由于 $T$ 自伴，$M_n$ 对 $T$ 不变：若 $x\in M_n$，则对每个 $k\le n$，
\[
\langle Tx,e_k\rangle=\langle x,Te_k\rangle=\lambda_k\langle x,e_k\rangle=0.
\]
因此限制算子 $T_n:=T|_{M_n}$ 仍是紧自伴算子。若 $T_n\neq 0$，重复上一步便在 $M_n$ 中得到单位特征向量 $e_{n+1}$，对应特征值 $\lambda_{n+1}$，并满足
\[
|\lambda_{n+1}|=\|T_n\|\le |\lambda_n|.
\]
若某一步 $T_n=0$，则递归终止。于是得到有限或可数的标准正交特征向量族 $\{e_n\}$。

若该族无限，则必须有 $\lambda_n\to 0$。否则存在 $\varepsilon>0$ 与子列 $(\lambda_{n_k})$ 使 $|\lambda_{n_k}|\ge \varepsilon$。但
\[
Te_{n_k}=\lambda_{n_k}e_{n_k}
\]
不可能含有收敛子列，因为
\[
\|Te_{n_k}-Te_{n_\ell}\|^2
=|\lambda_{n_k}|^2+|\lambda_{n_\ell}|^2
\ge 2\varepsilon^2,\qquad k\neq \ell,
\]
这与紧性矛盾。

设
\[
M:=\overline{\operatorname{span}\{e_n\}}.
\]
由于每个 $e_n$ 是特征向量，和上面同样的计算表明 $M^\perp$ 对 $T$ 不变。若 $T|_{M^\perp}\neq 0$，则按前述构造又能在 $M^\perp$ 中找到新的特征向量，与 $\{e_n\}$ 的极大性矛盾。故
\[
T|_{M^\perp}=0.
\]
任取 $x\in H$，由 Hilbert 空间中的正交分解可写成
\[
x=\sum_n \langle x,e_n\rangle e_n+z,\qquad z\in M^\perp.
\]
于是
\[
Tx=\sum_n \lambda_n \langle x,e_n\rangle e_n,
\]
因为 $Tz=0$。该级数按范数收敛：一方面 $(\langle x,e_n\rangle)\in \ell^2$，另一方面 $(\lambda_n)$ 有界且在无限情形下趋于零，因此部分和构成 Cauchy 列。

最后说明谱的刻画。对每个非零特征值，有限重性已由上一命题给出。反过来，若 $\lambda\neq 0$ 不属于特征值，则由于 $\lambda_n\to 0$ 且 $\lambda\neq \lambda_n$ 对所有 $n$ 都成立，存在常数
\[
\delta:=\inf\bigl(\,|\lambda|,\ |\lambda-\lambda_n|:n\ge 1\,\bigr)>0.
\]
对任意分解
\[
x=\sum_n a_n e_n+z,\qquad z\in M^\perp,
\]
定义
\[
R_\lambda x:=\sum_n \frac{a_n}{\lambda-\lambda_n}e_n+\frac1\lambda z.
\]
由于
\[
\left|\frac1{\lambda-\lambda_n}\right|\le \frac1\delta,
\qquad
\left|\frac1\lambda\right|\le \frac1\delta,
\]
该级数给出一个有界算子 $R_\lambda$。再由表示式直接计算，
\[
(\lambda I-T)R_\lambda x=R_\lambda(\lambda I-T)x=x.
\]
因此 $\lambda I-T$ 可逆。综上，非零谱恰由这些非零特征值组成；若谱无限，则唯一可能的聚点只能是 $0$。
\end{proof}

\paragraph{谱定理与二次型}

\begin{remark}
定理~\ref{thm:compact-selfadjoint-spectral} 说明，紧自伴算子可以像对称矩阵一样被正交分解。对优化而言，这意味着许多无穷维二次泛函
\[
q(x)=\frac12\langle Tx,x\rangle-\langle b,x\rangle
\]
仍能在特征方向上逐坐标分析。有限维的 Hessian 对角化、主成分分析和半正定矩阵分解，都只是这一结论的坍缩版本。
\end{remark}

\begin{example}[协方差算子的谱分解]\label{ex:covariance-operator}
在 $L^2$ 或 RKHS 环境中，协方差算子常常是正、紧、自伴的，因此可由定理~\ref{thm:compact-selfadjoint-spectral} 分解为
\[
Tx=\sum_n \lambda_n \langle x,e_n\rangle e_n,\qquad \lambda_n\ge 0.
\]
把这个例子放在这里，是为了把谱定理直接连接到后文的核方法、表示学习与低秩近似：主成分方向就是特征向量，能量大小就是特征值。
\end{example}

\subsection*{有限维对照}

在有限维 Hilbert 空间里，定理~\ref{thm:compact-selfadjoint-spectral} 退化为“实对称矩阵可正交对角化”。因此半正定规划中的特征值分解、二次型标准形和迹内积，都可以视为本节谱定理在矩阵层面的压缩版。不同之处只在于：无穷维里非零谱可能是可数无限的，而紧性保证它只能向 $0$ 聚集。

\subsection*{习题（待补）}

\begin{exercise}[紧自伴谱定理的矩阵版占位]\label{exr:spectral-matrix-placeholder}
补一题：把定理~\ref{thm:compact-selfadjoint-spectral} 限制到有限维空间，说明它如何退化为实对称矩阵的正交对角化。
\end{exercise}

\begin{exercise}[二次型桥接题占位]\label{exr:quadratic-form-placeholder}
补一题：设 $T\ge 0$ 为紧自伴算子，分析二次泛函 $q(x)=\frac12\langle Tx,x\rangle-\langle b,x\rangle$ 在特征基下的表达，并解释何时它退化为有限维半正定二次规划。
\end{exercise}

